\documentclass[a4paper,12pt]{article}
\usepackage[utf8]{inputenc}
\usepackage[ngerman]{babel}
\usepackage{geometry}
\geometry{a4paper, left=3cm, right=2.5cm, top=2.5cm, bottom=2.5cm}

\usepackage[hidelinks]{hyperref}
\usepackage{enumitem}
\usepackage{longtable}
\usepackage{array}
\usepackage{graphicx}
\usepackage{float}
\usepackage[table]{xcolor}
\usepackage[most]{tcolorbox}
\usepackage{tabularx}
\usepackage{titlesec}

% Neue Spaltendefinition für bessere Formatierung
\newcolumntype{L}[1]{>{\raggedright\arraybackslash}p{#1}}

\titleformat{\subsubsection}[runin]{\bfseries}{\thesubsubsection}{1em}{}[.]

\title{Testkonzept\\[1ex]\large{Projektname: \textit{TippHub by THM}}}
\author{Gruppe C}
\date{\today}

\begin{document}

\maketitle
\tableofcontents
\newpage

\section{Kommerzielle Wettanbieter (KWA)}

\subsection{Registrierung auf der Plattform (UC-24)}

\subsubsection{Testfall :UC-24-HS: Erfolgreiche Registrierung}
\begin{longtable}{|L{5cm}|L{10cm}|}
\hline
\textbf{Angaben für die Durchführung eines Tests} & Erfolgreiche Registrierung eines kommerziellen Wettanbieters \\
\hline
\textbf{Vorbedingungen} & Registrierungsformular ist geöffnet, Plattform ist betriebsbereit \\
\hline
\textbf{Nachbedingungen} & Wettanbieter ist erfolgreich registriert und wird zur Login-Seite weitergeleitet \\
\hline
\textbf{Eingabewerte} & 
\begin{minipage}[t]{\linewidth}
\vspace{0.5em}
\begin{itemize}[leftmargin=*]
  \item Firmenname: \texttt{SportwettenPro}
  \item E-Mail: \texttt{kontakt@sportwettenpro.de}
  \item Bankdaten: \textit{DE03 3455 4567 4567 45 BIC: bccbxä Bankname: ING-Bank}
  \item Passwort: \texttt{Sicher123!}
  \item Lizenznummer: \texttt{DE-123456}
\end{itemize}
\vspace{0.5em}
\end{minipage} \\
\hline
\textbf{Erwartete Ausgabewerte} & \texttt{„Registrierung erfolgreich“}, automatischer Wechsel zur Loginseite oder Dashboard \\
\hline
\textbf{Prüfanweisungen} &
\begin{minipage}[t]{\linewidth}
\vspace{0.5em}
\begin{itemize}[leftmargin=*]
  \item Felder ausfüllen
  \item Button „Registrieren“ drücken
  \item Meldung und Weiterleitung prüfen
  \item Datenbankeintrag kontrollieren
\end{itemize}
\vspace{0.5em}
\end{minipage} \\
\hline
\end{longtable}

\vspace{1.5cm}

\subsubsection{Testfall :UC-24-AS: Fehlgeschlagene Registrierung (Pflichtfeld leer)}
\begin{longtable}{|L{5cm}|L{10cm}|}
\hline
\textbf{Angaben für die Durchführung eines Tests} & Fehlgeschlagene Registrierung (Pflichtfeld leer) \\
\hline
\textbf{Vorbedingungen} & Registrierungsformular ist geöffnet \\
\hline
\textbf{Nachbedingungen} & Registrierung wird verweigert, Nutzer bleibt auf aktueller Seite \\
\hline
\textbf{Eingabewerte} &
\begin{minipage}[t]{\linewidth}
\vspace{0.5em}
\begin{itemize}[leftmargin=*]
  \item Firmenname: \texttt{SportwettenX}
  \item E-Mail: \textit{leer}
  \item Bankdaten: \textit{DE03 3455 4567 4567 45 BIC: bccbxä, ING-Bank}
  \item Passwort: \texttt{1234}
  \item Lizenznummer: \texttt{DE-999999}
\end{itemize}
\vspace{0.5em}
\end{minipage} \\
\hline
\textbf{Erwartete Ausgabewerte} & \texttt{„E-Mail darf nicht leer sein“}, keine Weiterleitung, keine Registrierungsmöglichkeit \\
\hline
\textbf{Prüfanweisungen} &
\begin{minipage}[t]{\linewidth}
\vspace{0.5em}
\begin{itemize}[leftmargin=*]
  \item Nur einige Felder ausfüllen und dabei E-Mail leer lassen
  \item „Registrieren“-Button drücken
  \item Fehlermeldung prüfen
  \item Sicherstellen, dass kein Nutzer erstellt wurde
\end{itemize}
\vspace{0.5em}
\end{minipage} \\
\hline
\end{longtable}





\subsection{Ankündigung für Wettangebote erstellen (UC-25)}

\subsubsection{Testfall: UC-25-HS: Ankündigung erfolgreich erstellen}
\begin{longtable}{|L{5cm}|L{10cm}|}
\hline
\textbf{Angaben für die Durchführung eines Tests} & Kommerzieller Wettanbieter (KWA) erstellt eine neue Ankündigung, um auf ein spezifisches Wettangebot aufmerksam zu machen \\
\hline
\textbf{Vorbedingungen} &
\begin{minipage}[t]{\linewidth}
\vspace{0.5em}
\begin{itemize}[leftmargin=*]
  \item KWA ist eingeloggt
  \item URL des Angebots ist freigegeben (Whitelist)
  \item Ausreichendes THMCoins-Guthaben vorhanden (groeser oder gleich 10)
\end{itemize}
\vspace{0.5em}
\end{minipage} \\
\hline
\textbf{Nachbedingungen} &
\begin{minipage}[t]{\linewidth}
\vspace{0.5em}
\begin{itemize}[leftmargin=*]
  \item Ankündigung wird erfolgreich veröffentlicht
  \item Meldung: „Ankündigung veröffentlicht“ erscheint
  \item Ankündigung ist im Dashboard sichtbar
\end{itemize}
\vspace{0.5em}
\end{minipage} \\
\hline
\textbf{Eingabewerte} &
\begin{minipage}[t]{\linewidth}
\vspace{0.5em}
\begin{itemize}[leftmargin=*]
  \item Titel: \texttt{„Top-Wette im Juli“}
  \item URL: \texttt{https://www.wettanbieter.de/juliaktion}
  \item Zielgruppe: Sportwettfreunde
\end{itemize}
\vspace{0.5em}
\end{minipage} \\
\hline
\textbf{Erwartete Ausgabewerte} &
\begin{minipage}[t]{\linewidth}
\vspace{0.5em}
\begin{itemize}[leftmargin=*]
  \item Meldung: \texttt{„Ankündigung veröffentlicht“}
  \item Eintrag im Ankündigungsmodul sichtbar
\end{itemize}
\vspace{0.5em}
\end{minipage} \\
\hline
\textbf{Prüfanweisungen} &
\begin{minipage}[t]{\linewidth}
\vspace{0.5em}
\begin{itemize}[leftmargin=*]
  \item Navigiere zu „Neue Ankündigung erstellen“
  \item Pflichtfelder ausfüllen und Formular absenden
  \item Bestätigungsnachricht prüfen
  \item Sichtbarkeit der Ankündigung im Dashboard verifizieren
\end{itemize}
\vspace{0.5em}
\end{minipage} \\
\hline
\end{longtable}

\vspace{1cm}

\subsubsection{Testfall: UC-25-AS-1.1: Ankündigung nicht möglich – URL nicht freigegeben}
\begin{longtable}{|L{5cm}|L{10cm}|}
\hline
\textbf{Angaben für die Durchführung eines Tests} & KWA versucht eine Ankündigung zu veröffentlichen – die URL ist jedoch nicht auf der Whitelist \\
\hline
\textbf{Vorbedingungen} &
\begin{minipage}[t]{\linewidth}
\vspace{0.5em}
\begin{itemize}[leftmargin=*]
  \item KWA ist eingeloggt
  \item Zugriff auf das Formular „Neue Ankündigung erstellen“ ist möglich
  \item THMCoins-Guthaben groesser oder gleich 10 vorhanden
\end{itemize}
\vspace{0.5em}
\end{minipage} \\
\hline
\textbf{Nachbedingungen} &
\begin{minipage}[t]{\linewidth}
\vspace{0.5em}
\begin{itemize}[leftmargin=*]
  \item Keine Veröffentlichung der Ankündigung
  \item Fehlermeldung wird angezeigt
\end{itemize}
\vspace{0.5em}
\end{minipage} \\
\hline
\textbf{Eingabewerte} &
\begin{minipage}[t]{\linewidth}
\vspace{0.5em}
\begin{itemize}[leftmargin=*]
  \item Titel: \texttt{„Sommerwette 2025“}
  \item URL: \texttt{https://unbekannte-domain.de}
  \item Zielgruppe: Laufkundschaft
\end{itemize}
\vspace{0.5em}
\end{minipage} \\
\hline
\textbf{Erwartete Ausgabewerte} &
\begin{minipage}[t]{\linewidth}
\vspace{0.5em}
\begin{itemize}[leftmargin=*]
  \item System zeigt Fehlermeldung: \texttt{„URL nicht freigegeben“}
  \item Formular bleibt erhalten, keine Weiterleitung
\end{itemize}
\vspace{0.5em}
\end{minipage} \\
\hline
\textbf{Prüfanweisungen} &
\begin{minipage}[t]{\linewidth}
\vspace{0.5em}
\begin{itemize}[leftmargin=*]
  \item Formular vollständig ausfüllen mit nicht freigegebener URL
  \item Absenden → Fehlermeldung prüfen
  \item Sicherstellen, dass keine Ankündigung erstellt wurde
\end{itemize}
\vspace{0.5em}
\end{minipage} \\
\hline
\end{longtable}

\vspace{1cm}

\subsubsection{Testfall: UC-27-AS-1.2: Guthaben kleiner als 10 THMCoins}
\begin{longtable}{|L{5cm}|L{10cm}|}
\hline
\textbf{Angaben für die Durchführung eines Tests} & KWA möchte eine Ankündigung veröffentlichen, scheitert jedoch wegen unzureichendem Guthaben \\
\hline
\textbf{Vorbedingungen} &
\begin{minipage}[t]{\linewidth}
\vspace{0.5em}
\begin{itemize}[leftmargin=*]
  \item KWA ist eingeloggt
  \item Zugriff auf das Formular „Neue Ankündigung erstellen“ ist möglich
  \item URL ist whitelisted
\end{itemize}
\vspace{0.5em}
\end{minipage} \\
\hline
\textbf{Nachbedingungen} &
\begin{minipage}[t]{\linewidth}
\vspace{0.5em}
\begin{itemize}[leftmargin=*]
  \item Ankündigung wird nicht veröffentlicht
  \item Zahlungsaufforderung erscheint
  \item Nutzer wird zur Guthabenaufladung weitergeleitet
\end{itemize}
\vspace{0.5em}
\end{minipage} \\
\hline
\textbf{Eingabewerte} &
\begin{minipage}[t]{\linewidth}
\vspace{0.5em}
\begin{itemize}[leftmargin=*]
  \item Titel: \texttt{„Superquote Bundesliga“}
  \item URL: \texttt{https://www.wettanbieter.de/superquote}
  \item Zielgruppe: Fußballfans
  \item Guthaben: \texttt{5 THMCoins}
\end{itemize}
\vspace{0.5em}
\end{minipage} \\
\hline
\textbf{Erwartete Ausgabewerte} &
\begin{minipage}[t]{\linewidth}
\vspace{0.5em}
\begin{itemize}[leftmargin=*]
  \item System zeigt Hinweis: \texttt{„Bitte laden Sie Ihr Guthaben auf“}
  \item Weiterleitung zum Zahlungsbereich erfolgt
\end{itemize}
\vspace{0.5em}
\end{minipage} \\
\hline
\textbf{Prüfanweisungen} &
\begin{minipage}[t]{\linewidth}
\vspace{0.5em}
\begin{itemize}[leftmargin=*]
  \item Formular korrekt ausfüllen bei zu geringem Guthaben
  \item Absenden → Hinweis und Weiterleitung prüfen
  \item Sicherstellen, dass keine Ankündigung erstellt wurde
\end{itemize}
\vspace{0.5em}
\end{minipage} \\
\hline
\end{longtable}


\subsection{Zielgruppe für Ankündigung festlegen}

\subsubsection{Testfall: UC-26-HS: Zielgruppe erfolgreich festlegen}
\begin{longtable}{|L{5cm}|L{10cm}|}
\hline
\textbf{Angaben für die Durchführung eines Tests} & Kommerzieller Wettanbieter (KWA) konfiguriert beim Erstellen einer Ankündigung eine Zielgruppe mittels Geofilter, Interessen und Aktivitätslevel. Das System prüft die Gültigkeit und meldet Rückgaben entsprechend \\
\hline
\textbf{Vorbedingungen} &
\begin{minipage}[t]{\linewidth}
\vspace{0.5em}
\begin{itemize}[leftmargin=*]
  \item KWA ist eingeloggt
  \item Erstellungsformular für Zielgruppen-Ankündigung ist geöffnet
\end{itemize}
\vspace{0.5em}
\end{minipage} \\
\hline
\textbf{Nachbedingungen} &
\begin{minipage}[t]{\linewidth}
\vspace{0.5em}
\begin{itemize}[leftmargin=*]
  \item Zielgruppe ist gespeichert
  \item Digitale Signatur wurde generiert
  \item Ankündigung ist validiert und bereit zur Verteilung
\end{itemize}
\vspace{0.5em}
\end{minipage} \\
\hline
\textbf{Eingabewerte} &
\begin{minipage}[t]{\linewidth}
\vspace{0.5em}
\begin{itemize}[leftmargin=*]
  \item Geofilter: Hessen, NRW
  \item Interessen: Fußball, Live-Wetten
  \item Aktivitätslevel: Hoch
  \item Link zur Aktion: \texttt{https://www.wettanbieter.de/topspiel}
\end{itemize}
\vspace{0.5em}
\end{minipage} \\
\hline
\textbf{Erwartete Ausgabewerte} &
\begin{minipage}[t]{\linewidth}
\vspace{0.5em}
\begin{itemize}[leftmargin=*]
  \item Meldung: \texttt{„Zielgruppe gespeichert“}
  \item System erstellt digitale Signatur
  \item Validierung erfolgreich
\end{itemize}
\vspace{0.5em}
\end{minipage} \\
\hline
\textbf{Prüfanweisungen} &
\begin{minipage}[t]{\linewidth}
\vspace{0.5em}
\begin{itemize}[leftmargin=*]
  \item Targeting-Optionen aufrufen und Zielgruppe konfigurieren
  \item Speichern → Rückmeldung und Validierung prüfen
  \item Nach Signatur-Erstellung Verteilung anstoßen
\end{itemize}
\vspace{0.5em}
\end{minipage} \\
\hline
\end{longtable}


\subsubsection{Testfall: UC-26-AS.1-2: Zielgruppenkonfiguration – Warnung oder Fehler}

\begin{longtable}{|L{5cm}|L{10cm}|}
\hline
\textbf{Angaben für die Durchführung eines Tests} & Zielgruppeneinstellungen führen zu Warn- oder Fehlermeldungen (z.B. wegen zu geringer Reichweite oder ungültiger Konfiguration) \\
\hline
\textbf{Vorbedingungen} &
\begin{minipage}[t]{\linewidth}
\vspace{0.3em}
\begin{itemize}[leftmargin=*]
  \item KWA ist eingeloggt
  \item Formular „Zielgruppen-Ankündigung“ ist geöffnet
\end{itemize}
\vspace{0.3em}
\end{minipage} \\
\hline
\textbf{Nachbedingungen} &
\begin{minipage}[t]{\linewidth}
\vspace{0.3em}
\begin{itemize}[leftmargin=*]
  \item Keine Speicherung oder Verteilung bei Fehlkonfiguration
  \item Bei Warnung darf gespeichert, aber nicht sofort verteilt werden
\end{itemize}
\vspace{0.3em}
\end{minipage} \\
\hline
\textbf{Eingabewerte – Variante 1 (Warnung)} &
\begin{minipage}[t]{\linewidth}
\vspace{0.3em}
\begin{itemize}[leftmargin=*]
  \item Region: Bremen
  \item Interessen: Pferdewetten
  \item Systemreichweite: < 5 Empfänger
\end{itemize}
\vspace{0.3em}
\end{minipage} \\
\hline
\textbf{Erwartete Ausgabewerte – Variante 1} &
\begin{minipage}[t]{\linewidth}
\vspace{0.3em}
\begin{itemize}[leftmargin=*]
  \item Warnmeldung: \texttt{„Geringe Reichweite“}
  \item Speicherung möglich, Verteilung optional
\end{itemize}
\vspace{0.3em}
\end{minipage} \\
\hline
\textbf{Eingabewerte – Variante 2 (Fehler)} &
\begin{minipage}[t]{\linewidth}
\vspace{0.3em}
\begin{itemize}[leftmargin=*]
  \item Ungültige Kombination von Parametern (z.B. leerer Filter)
\end{itemize}
\vspace{0.3em}
\end{minipage} \\
\hline
\textbf{Erwartete Ausgabewerte – Variante 2} &
\begin{minipage}[t]{\linewidth}
\vspace{0.3em}
\begin{itemize}[leftmargin=*]
  \item Fehlermeldung: \texttt{„Konfiguration ungültig“}
  \item Keine Speicherung möglich
\end{itemize}
\vspace{0.3em}
\end{minipage} \\
\hline
\textbf{Prüfanweisungen} &
\begin{minipage}[t]{\linewidth}
\vspace{0.3em}
\begin{itemize}[leftmargin=*]
  \item Variante 1: Zielgruppe absichtlich zu eng wählen → Warnung prüfen
  \item Variante 2: Pflichtfelder leer lassen oder ungültig kombinieren → Fehlermeldung prüfen
  \item Sicherstellen, dass keine Verteilung erfolgt
\end{itemize}
\vspace{0.3em}
\end{minipage} \\
\hline
\end{longtable}



\subsection{Via externem Link ankündigen}

\subsubsection{Testfall: UC-27-HS: Via externem Link die PUS ankündigen}
\vspace{0.5cm}

\begin{center}
\begin{tabular}{|p{5cm}|p{10.5cm}|}
\hline
\textbf{Angaben für die Durchführung eines Tests} & 
KWA nutzt einen externen Link, um die Promotion Unit (PUS) einer Community gezielt anzukündigen. \\
\hline
\textbf{Vorbedingungen} & 
\begin{minipage}[t]{10.5cm}
\begin{itemize}[leftmargin=*]
  \item KWA ist eingeloggt
  \item Ziel-Link ist freigegeben
\end{itemize}
\end{minipage} \\
\hline
\textbf{Nachbedingungen} & 
\begin{minipage}[t]{10.5cm}
\begin{itemize}[leftmargin=*]
  \item PUS wird erfolgreich beworben
  \item Ankündigung ist sichtbar
\end{itemize}
\end{minipage} \\
\hline
\textbf{Eingabewerte} & 
\begin{minipage}[t]{10.5cm}
\begin{itemize}[leftmargin=*]
  \item Externer Link zur PUS
  \item Zielgruppe: Community XYZ
\end{itemize}
\end{minipage} \\
\hline
\textbf{Erwartete Ausgabewerte} &
\begin{minipage}[t]{10.5cm}
\begin{itemize}[leftmargin=*]
  \item Meldung: \texttt{„Ankündigung veröffentlicht“}
  \item Verlinkung aufrufbar und gültig
\end{itemize}
\end{minipage} \\
\hline
\textbf{Prüfanweisungen} &
\begin{minipage}[t]{10.5cm}
\begin{itemize}[leftmargin=*]
  \item Externen Link eingeben
  \item Zielgruppe festlegen
  \item Sichtbarkeit und Gültigkeit prüfen
\end{itemize}
\end{minipage} \\
\hline
\end{tabular}
\end{center}

\vspace{1.5cm}



\subsubsection{Testfall: UC-27-AS – Fehlgeleitete Weiterleitung via externem Link}
\begin{tabular}{|p{5cm}|p{10.5cm}|}
\hline
\textbf{Angaben für die Durchführung eines Tests} & Der hinterlegte externe Link führt nicht zur PUS oder zeigt fehlerhafte Inhalte. \\
\hline
\textbf{Vorbedingungen} & 
\begin{itemize}[leftmargin=*]
  \item KWA hat Ankündigung mit fehlerhaftem Link erstellt
\end{itemize} \\
\hline
\textbf{Nachbedingungen} & 
\begin{itemize}[leftmargin=*]
  \item Weiterleitung scheitert oder führt ins Leere
\end{itemize} \\
\hline
\textbf{Eingabewerte} &
\begin{itemize}[leftmargin=*]
  \item Falscher Link zur PUS
\end{itemize} \\
\hline
\textbf{Erwartete Ausgabewerte} & 
\begin{itemize}[leftmargin=*]
  \item Fehlermeldung oder leere Zielseite
  \item Eventuell Logging des Fehlers
\end{itemize} \\
\hline
\textbf{Prüfanweisungen} &
\begin{itemize}[leftmargin=*]
  \item Ankündigung mit ungültigem Link öffnen
  \item Weiterleitung prüfen
\end{itemize} \\
\hline
\end{tabular}

\subsection{Durch externen Link weiterleiten}

\subsubsection{Testfall: UC-28-HS – Externer Link funktioniert}
\vspace{0.5cm}

\begin{center}
\begin{tabular}{|p{5cm}|p{10.5cm}|}
\hline
\textbf{Angaben für die Durchführung eines Tests} & Nutzer klickt auf externen Link in der Ankündigung und wird korrekt zur Zielseite weitergeleitet. \\
\hline
\textbf{Vorbedingungen} &
\begin{minipage}[t]{10.5cm}
\begin{itemize}[leftmargin=*]
  \item Externer Link wurde korrekt angegeben
  \item Zielseite ist erreichbar
\end{itemize}
\end{minipage} \\
\hline
\textbf{Nachbedingungen} &
\begin{minipage}[t]{10.5cm}
\begin{itemize}[leftmargin=*]
  \item Nutzer sieht gewünschte Zielseite
\end{itemize}
\end{minipage} \\
\hline
\textbf{Eingabewerte} & 
\begin{minipage}[t]{10.5cm}
\begin{itemize}[leftmargin=*]
  \item Link: \texttt{https://partnerseite.de/aktion}
\end{itemize}
\end{minipage} \\
\hline
\textbf{Erwartete Ausgabewerte} &
\begin{minipage}[t]{10.5cm}
\begin{itemize}[leftmargin=*]
  \item Zielseite wird geöffnet
\end{itemize}
\end{minipage} \\
\hline
\textbf{Prüfanweisungen} &
\begin{minipage}[t]{10.5cm}
\begin{itemize}[leftmargin=*]
  \item Auf Link in Ankündigung klicken
  \item Prüfen, ob richtige Seite erscheint
\end{itemize}
\end{minipage} \\
\hline
\end{tabular}
\end{center}

\vspace{1.5cm}



\subsubsection{Testfall: UC-28-AS – Externer Link nicht erlaubt}
\vspace{0.5cm}

\begin{center}
\begin{tabular}{|p{5cm}|p{10.5cm}|}
\hline
\textbf{Angaben für die Durchführung eines Tests} & 
System erkennt, dass ein angegebener externer Link nicht auf der Whitelist steht. \\
\hline
\textbf{Vorbedingungen} &
\begin{minipage}[t]{10.5cm}
\begin{itemize}[leftmargin=*]
  \item Ankündigung mit nicht genehmigtem Link wird erstellt
\end{itemize}
\end{minipage} \\
\hline
\textbf{Nachbedingungen} &
\begin{minipage}[t]{10.5cm}
\begin{itemize}[leftmargin=*]
  \item Link wird nicht gespeichert
  \item Hinweis auf nicht erlaubten Link erscheint
\end{itemize}
\end{minipage} \\
\hline
\textbf{Eingabewerte} & 
\begin{minipage}[t]{10.5cm}
\begin{itemize}[leftmargin=*]
  \item Link: \texttt{https://nicht-erlaubt.com}
\end{itemize}
\end{minipage} \\
\hline
\textbf{Erwartete Ausgabewerte} &
\begin{minipage}[t]{10.5cm}
\begin{itemize}[leftmargin=*]
  \item Fehlermeldung: \texttt{„Externer Link nicht erlaubt“}
\end{itemize}
\end{minipage} \\
\hline
\textbf{Prüfanweisungen} &
\begin{minipage}[t]{10.5cm}
\begin{itemize}[leftmargin=*]
  \item Link einfügen
  \item Meldung prüfen
\end{itemize}
\end{minipage} \\
\hline
\end{tabular}
\end{center}

\vspace{1.5cm}



\subsection{THMCoins an PUS gutschreiben (UC-29)}

\subsubsection{Testfall: UC-29-HS – THMCoins erfolgreich gutschreiben}
\vspace{0.5cm}

\renewcommand{\arraystretch}{1.3}
\begin{center}
\begin{tabular}{|p{5cm}|p{10.5cm}|}
\hline
\textbf{Angaben für die Durchführung eines Tests} & 
PUS erhält eine definierte Anzahl an THMCoins von KWA als Vergütung. \\
\hline
\textbf{Vorbedingungen} &
\begin{minipage}[t]{10.5cm}
\begin{itemize}[leftmargin=*]
  \item KWA hat ausreichendes Guthaben
  \item PUS ist bekannt und aktiv
\end{itemize}
\vspace{0.2cm}
\end{minipage} \\
\hline
\textbf{Nachbedingungen} &
\begin{minipage}[t]{10.5cm}
\begin{itemize}[leftmargin=*]
  \item Coins werden korrekt übertragen
  \item Transaktion wird dokumentiert
\end{itemize}
\vspace{0.2cm}
\end{minipage} \\
\hline
\textbf{Eingabewerte} &
\begin{minipage}[t]{10.5cm}
\begin{itemize}[leftmargin=*]
  \item PUS-ID: \texttt{PUS-7215}
  \item Anzahl: \texttt{15 THMCoins}
\end{itemize}
\vspace{0.2cm}
\end{minipage} \\
\hline
\textbf{Erwartete Ausgabewerte} &
\begin{minipage}[t]{10.5cm}
\begin{itemize}[leftmargin=*]
  \item Buchungseintrag in beiden Konten
  \item Erfolgsmeldung
\end{itemize}
\vspace{0.2cm}
\end{minipage} \\
\hline
\textbf{Prüfanweisungen} &
\begin{minipage}[t]{10.5cm}
\begin{itemize}[leftmargin=*]
  \item Transaktion starten
  \item Gutschrift prüfen
\end{itemize}
\vspace{0.2cm}
\end{minipage} \\
\hline
\end{tabular}
\end{center}

\vspace{1.5cm}

\subsubsection{Testfall: UC-29-AS – Gutschrift nicht möglich: unzureichendes Guthaben}
\vspace{0.5cm}

\renewcommand{\arraystretch}{1.3}
\begin{center}
\begin{tabular}{|p{5cm}|p{10.5cm}|}
\hline
\textbf{Angaben für die Durchführung eines Tests} & 
Versuch, einer PUS THMCoins gutzuschreiben, obwohl das KWA-Guthaben nicht ausreicht. \\
\hline
\textbf{Vorbedingungen} &
\begin{minipage}[t]{10.5cm}
\begin{itemize}[leftmargin=*]
  \item KWA ist angemeldet
  \item PUS-ID ist gültig
  \item KWA-Guthaben < angeforderter Betrag
\end{itemize}
\vspace{0.2cm}
\end{minipage} \\
\hline
\textbf{Nachbedingungen} &
\begin{minipage}[t]{10.5cm}
\begin{itemize}[leftmargin=*]
  \item Keine Gutschrift erfolgt
  \item Hinweis auf fehlendes Guthaben wird angezeigt
\end{itemize}
\vspace{0.2cm}
\end{minipage} \\
\hline
\textbf{Eingabewerte} &
\begin{minipage}[t]{10.5cm}
\begin{itemize}[leftmargin=*]
  \item PUS-ID: \texttt{PUS-7215}
  \item Betrag: \texttt{15 THMCoins}
\end{itemize}
\vspace{0.2cm}
\end{minipage} \\
\hline
\textbf{Erwartete Ausgabewerte} &
\begin{minipage}[t]{10.5cm}
\begin{itemize}[leftmargin=*]
  \item Fehlermeldung: \texttt{„Guthaben nicht ausreichend“}
  \item Keine Buchung in Transaktionshistorien
\end{itemize}
\vspace{0.2cm}
\end{minipage} \\
\hline
\textbf{Prüfanweisungen} &
\begin{minipage}[t]{10.5cm}
\begin{itemize}[leftmargin=*]
  \item KWA mit wenig Guthaben anmelden
  \item Betrag zur Gutschrift festlegen
  \item Übertragung auslösen
  \item Meldung und Transaktionsverlauf kontrollieren
\end{itemize}
\vspace{0.2cm}
\end{minipage} \\
\hline
\end{tabular}
\end{center}

\vspace{1.5cm}

\section*{Testfall: Monatliche THMCoins-Gutschrift an KWA gutschreiben}

\begin{tabular}{|p{5cm}|p{10.5cm}|}
\hline
\textbf{Angaben für die Durchführung eines Tests} & Monatliche Gutschrift der gesammelten THMCoins eines KWAs auf dessen internes Konto \\
\hline
\textbf{Vorbedingungen} & 
\begin{itemize}[leftmargin=*]
  \item Gültiger KWA-Account existiert
  \item Es liegen Transaktionen mit THMCoins vor
  \item Monatlicher Abrechnungszyklus ist aktiv
\end{itemize} \\
\hline
\textbf{Nachbedingungen} & 
\begin{itemize}[leftmargin=*]
  \item Gutschrift der Summe aller verdienten THMCoins im aktuellen Monat
  \item Buchung ist im KWA-Dashboard ersichtlich
  \item Abrechnungsbericht wird gespeichert
\end{itemize} \\
\hline
\textbf{Eingabewerte} &
\begin{itemize}[leftmargin=*]
  \item KWA-ID: \texttt{KWA-87912}
  \item Zeitraum: 01.06.2025 – 30.06.2025
  \item Transaktionen: 25 THMCoins, 35 THMCoins, 40 THMCoins
\end{itemize} \\
\hline
\textbf{Erwartete Ausgabewerte} & 
\begin{itemize}[leftmargin=*]
  \item Kontogutschrift: \texttt{100 THMCoins}
  \item Meldung im auf Display: \texttt{„THMCoins erfolgreich gutgeschrieben“}
\end{itemize} \\
\hline
\textbf{Prüfanweisungen} &
\begin{itemize}[leftmargin=*]
  \item Sicherstellen, dass automatische Abrechnung gestartet wird
  \item Kontostand vor und nach Ausführung dokumentieren
  \item Darstellung im Dashboard verifizieren
\end{itemize} \\
\hline
\end{tabular}

\end{document}